\code{IRPC} represents an Inferior Remote Procedure Call that can be run in a
target process. See \code{Process::postIRPC} and \code{Thread::postIRPC} for
information about posting an IRPC.

\begin{apient}
IRPC::ptr IRPC::const_ptr
\end{apient}

\apidesc{
The IRPC::ptr and IRPC::const\_ptr respectively represent a pointer and a const
pointer to an IRPC object. Both pointer types are reference counted and will
cause the underlying IRPC object to be cleaned when there are no more
references. ProcControlAPI will maintain internal references to any IRPC
currently posted or executing.
}

\begin{apient}
IRPC::ptr createIRPC(void *binary_blob, unsigned int size, bool non_blocking = false)
\end{apient}

\apidesc{
}

\begin{apient}
IRPC::ptr createIRPC(void *binary_blob, unsigned int size, Dyninst::Address addr, bool non_blocking = false)
\end{apient}

\apidesc{
The createIRPC static function creates and returns a new IRPC object. The
binary\_blob and size parameters specify a buffer of machine code bytes that
this IRPC should execute when invoked. ProcControlAPI will maintain its own
copy of the binary\_blob buffer, the ProcControlAPI user can free the buffer
once this function completes. 
If the non\_blocking parameter is true then calls to
Process::handleEvents will block until this IRPC is
completed. 
If the addr parameter is given, then ProcControlAPI will write and run the
binary code at addr. Otherwise ProcControlAPI will select a location at
which to run the IRPC.
}

\begin{apient}
Dyninst::Address getAddress() const
\end{apient}

\apidesc{
The getAddress function returns the address at which the IRPC will be run. If
the IRPC was not given an address at construction and has not yet started
running, then this function may return 0.
}

\begin{apient}
void *getBinaryCodeBlob() const
\end{apient}

\apidesc{
The getBinaryCodeBlob returns a pointer to memory that contains the binary code
for this IRPC. 
}

\begin{apient}
unsigned int getBinaryCodeSize() const
\end{apient}

\apidesc{
The getBinaryCodeSize function returns the size of the binary code blob buffer.
}

\begin{apient}
unsigned long getID() const
\end{apient}

\apidesc{
The getID function returns an integer identifier that uniquely identifies this
IRPC.
}

\begin{apient}
void setStartOffset(unsigned long off)
\end{apient}

\apidesc{
By default an IRPC will start executing its code blob at the beginning of the
blob. This function can be used to tell ProcControlAPI to start execution of
the code blob at some byte offset, off, into the blob. 
This function should be called before the IRPC is posted.
}

\begin{apient}
unsigned long getStartOffset() const
\end{apient}

\apidesc{
If a start offset has been set for this IRPC, then getStartOffset returns it.
Otherwise this function returns 0.
}
